\begin{lemma}\label{lem:decic333-finite-image}
Put $F=\mathbb Q(\zeta_3)$ and $X=X_1(3,33)$, and let
$h(E,P_3,Q_{33})=(E,P_3,10Q_{33})$.
Then $(h_*-1)\operatorname{Jac}(X)(F)$ is finite.
\end{lemma}
\begin{proof}
Over $F$, the curve $X$ is isomorphic to the base change of
$X_H(99)$, where $H=\langle34\rangle$; see \cite[Section~2]{DS}.
The involution $h$ is induced by the diamond operator
$\langle10\rangle$. The source has genus $81$ and its quotient
by $h$ has genus $41$, so its anti-invariant Jacobian factor
has dimension $40$.

The even nebentypes $\varepsilon$ modulo $99$ satisfying
$\varepsilon(34)=1$ and $\varepsilon(10)=-1$ have conductor
dividing $33$. They have precisely two Galois orbits, of orders
$2$ and $10$. In each case the corresponding cuspidal modular
symbol space with sign $+$ has dimension $8$ over the character
field. Their rational dimensions are $8$ and $4\cdot8$, giving
all $40$ dimensions of the anti-invariant factor.
The exact winding computation gives full Hecke orbit rank
$8$ in both spaces. This is certified by reduction modulo $5$
and $11$, respectively, avoiding the finitely many denominator
primes. Thus every newform contributing to these spaces has
nonzero central $L$-value.

Let $\chi_{-3}$ be the quadratic character of $F$.
The twist-containment argument at the end of the proof of
\cite[Theorem~4.1]{DS}, using Atkin--Li, shows that twisting a
form in $S_2(\Gamma_0(99),\varepsilon)$ by $\chi_{-3}$ stays in
this same space: the conductor of $\varepsilon$ divides $33$
and $99=3\cdot33$. In particular its nebentype and its $h$-sign
are unchanged. The preceding winding computation therefore
also covers every $\chi_{-3}$-twist of an anti-invariant factor.
The nonvanishing theorem for modular abelian varieties used
in \cite[Section~4]{DS} gives rank zero over $\mathbb Q$ for
each factor and its twist, hence rank zero over $F$.
The image of $h_*-1$ is contained in this factor, proving the
claim. The winding vectors and Hecke matrices are retained,
and their full ranks are checked independently in Sage.
\end{proof}
